\chapter{周秦卷索引}

本索引只收录本卷已经写入正文、且可由链接抵达的事件，以及已经置入正文的图件；它不是事件账本的替代品。年代有争议的节点仍依正文的说明处理，未写成正文的账本节点不在此增列。

\section{中国事件索引}

\subsection*{西周}
\begin{itemize}
  \item \eventref{WZ-1046-MUYE}{西周·牧野}；\eventref{WZ-1040-DONGZHENG}{西周·三监东征}；\eventref{WZ-EAST-CHENGZHOU}{西周·营建成周}；\eventref{WZ-EAST-ENFEOFFMENT}{西周·重组东方封国网络}。
  \item \eventref{WZ-CHENGKANG-CONSOLIDATION}{西周·成康巩固}；\eventref{WZ-MU-TRAVELS}{西周·穆王远行}；\eventref{WZ-0977-ZHAONAN}{西周·昭王南征}。
  \item \eventref{WZ-0841-GONGHE}{西周·共和}；\eventref{WZ-0828-XUAN-ACCESSION}{西周·宣王即位}；\eventref{WZ-XUAN-REVIVAL}{西周·宣王中兴}；\eventref{WZ-0771-HAOJING}{西周·镐京陷落}。
\end{itemize}

\subsection*{春秋}
\begin{itemize}
  \item \eventref{SA-0770-EAST}{春秋·东迁}；\eventref{SA-0770-QIN-ENFEOFF}{春秋·秦获封关中}；\eventref{SA-0722-CHUNQIU}{春秋·鲁隐公元年}；\eventref{SA-0707-XUGE}{春秋·繻葛之战}；\eventref{SA-0704-CHU-KING}{春秋·楚君称王}。
  \item \eventref{SA-0685-QIHUAN}{春秋·齐桓公即位}；\eventref{SA-0656-SHAOLING}{春秋·召陵之盟}；\eventref{SA-0651-KUIQIU}{春秋·葵丘之会}；\eventref{SA-0636-JIN-WEN}{春秋·晋文公即位}；\eventref{SA-0632-CHENGPU}{春秋·城濮之战}；\eventref{SA-0627-XIAO}{春秋·殽之战}。
  \item \eventref{SA-0597-BI}{春秋·邲之战}；\eventref{SA-0575-YANLING}{春秋·鄢陵之战}；\eventref{SA-0546-MIBING}{春秋·弭兵之会}；\eventref{SA-0506-BOJU}{春秋·柏举之战}；\eventref{SA-0494-FUJIAO}{春秋·夫椒之战}。
  \item \eventref{SA-0484-WUZIXU-DEATH}{春秋·伍子胥之死}；\eventref{SA-0482-HUANGCHI}{春秋·黄池之会}；\eventref{SA-0481-LIN}{春秋·获麟与陈氏政变}；\eventref{SA-0481-END}{春秋·记事终点}；\eventref{SA-0473-WU-FALL}{春秋·越灭吴}。
\end{itemize}

\subsection*{战国与秦}
\begin{itemize}
  \item \eventref{WQ-0453-ZHI}{战国·三家灭智}；\eventref{WQ-0408-WEI}{战国·魏国早期改革}；\eventref{WQ-0403-THREEJIN}{战国·三家分晋}；\eventref{WQ-0386-TIANQI}{战国·田氏代齐}；\eventref{WQ-0381-WUQI-CHU}{战国·吴起变法}；\eventref{WQ-0375-HAN}{战国·韩灭郑}。
  \item \eventref{WQ-0362-WEIQIN}{战国·少梁之战}；\eventref{WQ-0356-SHANGYANG}{战国·商鞅变法}；\eventref{WQ-0354-SHENBUHAI}{战国·申不害相韩}；\eventref{WQ-0354-GUILIN}{战国·桂陵之战}；\eventref{WQ-0342-MALING}{战国·马陵之战}；\eventref{WQ-0334-KINGS}{战国·齐魏称王}。
  \item \eventref{WQ-0317-YAN}{战国·子之之乱}；\eventref{WQ-0314-QI-YAN}{战国·齐伐燕}；\eventref{WQ-0307-HUFU}{战国·胡服骑射}；\eventref{WQ-0307-QIN-SUCCESSION}{秦·昭襄王即位}；\eventref{WQ-0295-SHAQIU}{战国·沙丘之乱}；\eventref{WQ-0293-YIQUE}{战国·伊阙之战}。
  \item \eventref{WQ-0284-QI}{战国·燕破齐}；\eventref{WQ-0279-TIANDAN}{战国·田单复齐}；\eventref{WQ-0266-FANJU}{秦·范雎入秦与远交近攻}；\eventref{WQ-0262-SHANGDANG}{战国·上党归赵}；\eventref{WQ-0260-CHANGPING}{战国·长平之战}。
  \item \eventref{WQ-0256-ZHOUDIE}{战国·东周亡}；\eventref{WQ-0249-EASTZHOU}{战国·东周君国亡}；\eventref{QIN-0230-UNIFICATION}{秦·统一六国}；\eventref{QIN-0226-YAN}{秦·破蓟}；\eventref{QIN-0222-YAN}{秦·灭燕}。
  \item \eventref{QIN-0213-EMPIRE}{秦·高强度治理}；\eventref{QIN-0210-SUCCESSION}{秦·沙丘继承}；\eventref{QIN-0209-FALL}{秦·起事与覆亡}、\eventref{QIN-0209-CHENSHENG}{秦·大泽起事}；\eventref{QIN-0208-ZHANGHAN}{秦·章邯出关}；\eventref{QIN-0207-LISI}{秦·李斯被杀}；\eventref{QIN-0207-ZIYING}{秦·子婴即位}；\eventref{QIN-0207-JULU}{秦·巨鹿之战}；\eventref{QIN-0206-SURRENDER}{秦·子婴降汉}。
\end{itemize}

\section{政体与国家索引}

\begin{description}
  \item[周王室与西周] 牧野后的王室建构、东方经营与晚期危机：\eventref{WZ-1046-MUYE}{西周·牧野}、\eventref{WZ-EAST-CHENGZHOU}{西周·营建成周}、\eventref{WZ-EAST-ENFEOFFMENT}{西周·重组东方封国网络}、\eventref{WZ-0841-GONGHE}{西周·共和}、\eventref{WZ-0771-HAOJING}{西周·镐京陷落}。
  \item[东周王室] 东迁以后保有册命与名分，至战国退出政治舞台：\eventref{SA-0770-EAST}{春秋·东迁}、\eventref{SA-0707-XUGE}{春秋·繻葛之战}、\eventref{WQ-0403-THREEJIN}{战国·三家分晋}、\eventref{WQ-0256-ZHOUDIE}{战国·东周亡}、\eventref{WQ-0249-EASTZHOU}{战国·东周君国亡}。
  \item[齐] 齐桓与会盟、田氏易姓、伐燕及复齐：\eventref{SA-0685-QIHUAN}{春秋·齐桓公即位}、\eventref{SA-0656-SHAOLING}{春秋·召陵之盟}、\eventref{WQ-0386-TIANQI}{战国·田氏代齐}、\eventref{WQ-0314-QI-YAN}{战国·齐伐燕}、\eventref{WQ-0284-QI}{战国·燕破齐}、\eventref{WQ-0279-TIANDAN}{战国·田单复齐}。
  \item[晋、赵、魏、韩] 晋的霸权与卿族裂解，及其后的三国竞争：\eventref{SA-0636-JIN-WEN}{春秋·晋文公即位}、\eventref{SA-0632-CHENGPU}{春秋·城濮之战}、\eventref{SA-0575-YANLING}{春秋·鄢陵之战}、\eventref{WQ-0453-ZHI}{战国·三家灭智}、\eventref{WQ-0403-THREEJIN}{战国·三家分晋}；另见\eventref{WQ-0408-WEI}{战国·魏国早期改革}、\eventref{WQ-0375-HAN}{战国·韩灭郑}、\eventref{WQ-0307-HUFU}{战国·胡服骑射}与\eventref{WQ-0260-CHANGPING}{战国·长平之战}。
  \item[楚] 从自称王到江汉大国的北上与制度压力：\eventref{SA-0704-CHU-KING}{春秋·楚君称王}、\eventref{SA-0597-BI}{春秋·邲之战}、\eventref{SA-0506-BOJU}{春秋·柏举之战}、\eventref{WQ-0381-WUQI-CHU}{战国·吴起变法}。
  \item[秦] 关中边缘、变法与统一帝国：\eventref{SA-0770-QIN-ENFEOFF}{春秋·秦获封关中}、\eventref{SA-0627-XIAO}{春秋·殽之战}、\eventref{WQ-0356-SHANGYANG}{战国·商鞅变法}、\eventref{WQ-0350-QINREFORM}{战国·秦迁咸阳与第二次变法}、\eventref{QIN-0230-UNIFICATION}{秦·统一六国}、\eventref{QIN-0206-SURRENDER}{秦·子婴降汉}。
  \item[鲁、郑、宋] 编年书写、中原要冲与弭兵调停：\eventref{SA-0722-CHUNQIU}{春秋·鲁隐公元年}、\eventref{SA-0707-XUGE}{春秋·繻葛之战}、\eventref{SA-0546-MIBING}{春秋·弭兵之会}、\eventref{SA-0481-LIN}{春秋·获麟与陈氏政变}。
  \item[吴、越] 下游水网中的北上争盟与相互吞并：\eventref{SA-0506-BOJU}{春秋·柏举之战}、\eventref{SA-0494-FUJIAO}{春秋·夫椒之战}、\eventref{SA-0482-HUANGCHI}{春秋·黄池之会}、\eventref{SA-0473-WU-FALL}{春秋·越灭吴}。
  \item[燕] 内乱、复仇式动员与秦军东进：\eventref{WQ-0317-YAN}{战国·子之之乱}、\eventref{WQ-YAN-ZHAO}{战国·燕昭王蓄力}、\eventref{WQ-0284-QI}{战国·燕破齐}、\eventref{QIN-0226-YAN}{秦·破蓟}、\eventref{QIN-0222-YAN}{秦·灭燕}。
\end{description}

\section{主题索引}

\begin{description}
  \item[战争] 从王朝奠基的\eventref{WZ-1046-MUYE}{西周·牧野}，到春秋的\eventref{SA-0632-CHENGPU}{春秋·城濮之战}、\eventref{SA-0597-BI}{春秋·邲之战}、\eventref{SA-0506-BOJU}{春秋·柏举之战}，再到战国的\eventref{WQ-0354-GUILIN}{战国·桂陵之战}、\eventref{WQ-0342-MALING}{战国·马陵之战}、\eventref{WQ-0293-YIQUE}{战国·伊阙之战}与\eventref{WQ-0260-CHANGPING}{战国·长平之战}。参见秦末的\eventref{QIN-0207-JULU}{秦·巨鹿之战}。
  \item[东方经营与封国网络] 从\eventref{WZ-1040-DONGZHENG}{西周·三监东征}后的安置，到\eventref{WZ-EAST-CHENGZHOU}{西周·营建成周}与\eventref{WZ-EAST-ENFEOFFMENT}{西周·重组东方封国网络}，可见王室如何以成周、册命、驻守、祭祀和封国维系东方；这是一套须不断确认的关系网络，而非一次划定的行政区划。
  \item[改革与国家能力] \eventref{WQ-0408-WEI}{战国·魏国早期改革}、\eventref{WQ-WEI-WUQI}{战国·魏武卒}、\eventref{WQ-0381-WUQI-CHU}{战国·吴起变法}、\eventref{WQ-0354-SHENBUHAI}{战国·申不害相韩}、\eventref{WQ-0356-SHANGYANG}{战国·商鞅变法}与\eventref{WQ-0350-QINREFORM}{战国·秦迁咸阳与第二次变法}。另见图“战国变法关系示意”。
  \item[经济、动员与交通] 西周的成周和东方网络以\eventref{WZ-EAST-CHENGZHOU}{西周·营建成周}为入口；战国则可从\eventref{WQ-0350-QINREFORM}{战国·秦迁咸阳与第二次变法}、\eventref{WQ-0262-SHANGDANG}{战国·上党归赵}与\eventref{WQ-0260-CHANGPING}{战国·长平之战}观察田界、县邑、道路与补给怎样进入竞争。秦帝国的日常征发见\eventref{QIN-0213-EMPIRE}{秦·高强度治理}。
  \item[文化、礼制与书写] 西周王权的礼制语言见\eventref{WZ-CHENGKANG-CONSOLIDATION}{西周·成康巩固}；春秋编年视角的起点和终点见\eventref{SA-0722-CHUNQIU}{春秋·鲁隐公元年}、\eventref{SA-0481-END}{春秋·记事终点}。齐的稷下、游说与法家讨论置于战国“文化与思想”一节，并与\eventref{WQ-0386-TIANQI}{战国·田氏代齐}所后的齐国路径并读。
  \item[边疆、地理与迁徙] 南向通道的限度见\eventref{WZ-0977-ZHAONAN}{西周·昭王南征}，西北压力见\eventref{WZ-XUAN-REVIVAL}{西周·宣王中兴}；秦在关中的位置见\eventref{SA-0770-QIN-ENFEOFF}{春秋·秦获封关中}；赵的北方应变见\eventref{WQ-0307-HUFU}{战国·胡服骑射}；燕的辽东退守与覆亡见\eventref{QIN-0226-YAN}{秦·破蓟}、\eventref{QIN-0222-YAN}{秦·灭燕}。
  \item[史料与争议] 年代与材料层次的提示集中见\eventref{WZ-1046-MUYE}{西周·牧野}、\eventref{WZ-MU-TRAVELS}{西周·穆王远行}、\eventref{WZ-0841-GONGHE}{西周·共和}、\eventref{SA-0481-END}{春秋·记事终点}、\eventref{WQ-0279-TIANDAN}{战国·田单复齐}与\eventref{QIN-0209-CHENSHENG}{秦·大泽起事}。相关正文区分金文、传世编年与后出叙事的证据范围，不把后世故事当作现场记录。
\end{description}

\section{跨期制度得失：从封国网络到帝国行政}

这一卷所说的“得失”，不是给某种制度判定高下，也不是把后来的郡县制倒投回西周。比较的尺度是具体而有限的：在各自的地理、战争与资源条件下，它能否把人力和物资组织起来，能否让地方仍有可协商的空间，又把怎样的负担、风险与后续难题留给了谁。四个阶段并非一条预定通向帝国的直线；每一次重组，都是对上一阶段难题的一种回应，也会制造新的选择。

\begin{textwindow}[《诗经》里的“维新”]
“周虽旧邦，其命维新。”

语出《诗经·大雅·文王》。这首传世诗篇以周的受命为主题，不是一道可供逐条执行的分封诰命，也不能用来证明周初某项具体制度；它在这里所提示的，是政治秩序常以“受命”与“更新”的语言解释自身。以下的比较仍以各事件的金文、传世叙事、简牍与后出史书的证据边界为准。
\end{textwindow}

\subsection{西周：把空间交给关系网络}

从\eventref{WZ-1040-DONGZHENG}{西周·三监东征}到\eventref{WZ-EAST-CHENGZHOU}{西周·营建成周}，王室面对的是征服之后如何经营东方的问题。封国、宗族、册命与祭祀并非一张整齐的行政区划图，却能让周王不必在每一处常驻同等规模的官吏和军队；地方精英以本地的武装、劳作与人际关系承担守卫、开拓和日常调解，王室则以名分、赏赐、会集与征伐维持网络。它的收益是以较少的直接驻守覆盖较大的空间，成本则是筑城、迁置、供给与服役须由王室、封国及其辖下人群共同负担，而现存材料很少留下后者自己的账目。

这种安排的要害在于不断兑现关系，而不在一次授土后即可永久命令。地方获得承认，也随之积累代际的资源、武力与声望；王室强盛时，这些能力是网络的支点，王室失去调解和供给能力时，它们又成为各方自行结盟的条件。\eventref{WZ-0771-HAOJING}{西周·镐京陷落}与其后的东迁并不证明封国安排必然导致覆亡，却清楚显示：礼仪名分可以保存，支撑它的军事与资源关系却未必能一同保存。

\subsection{春秋：以霸主调度，而不把诸侯变成属县}

\eventref{SA-0770-EAST}{春秋·东迁}以后，周王室保留的名分已不足以独自协调诸侯。齐桓在\eventref{SA-0685-QIHUAN}{春秋·齐桓公即位}后组织会盟，晋又在\eventref{SA-0632-CHENGPU}{春秋·城濮之战}后成为中原枢纽；霸主以兵力、资源和“尊王”的语言，为救援、边患与争端提供临时中心。与长期占领相比，这是一种成本较低的秩序：小国仍可借会盟申诉或制衡，大国也不必立刻承担每一国的赋税、继承和地方治理。

代价同样来自它的弹性。盟约没有常设官署和稳定的执行链，霸主的继承危机、卿族的分权或外部战争都能迫使各国重新计算利害。\eventref{SA-0546-MIBING}{春秋·弭兵之会}能够暂缓冲突，恰好说明调停仍需各方当场接受，不能自行变成日常政府。春秋的霸政因而不是“失败的统一”，而是在多中心世界里降低协调成本的一种办法；它留下的难题是，如何在不完全依靠贵族关系与临时承认的条件下，稳定取得兵员、粮食和地方信息。

\subsection{战国：把可调用的资源写入县邑、户籍与赏罚}

\eventref{WQ-0403-THREEJIN}{战国·三家分晋}使卿族掌握国家资源的后果变得格外可见。魏、楚、韩、秦的改革路径并不相同，却都在竞争中尝试使土地、人口、军队与官职更可登记、考核和调度。以\eventref{WQ-0356-SHANGYANG}{战国·商鞅变法}和\eventref{WQ-0350-QINREFORM}{战国·秦迁咸阳与第二次变法}为例，军功、编户、连带责任、县邑、田界和度量等安排，缩短了统治者从战争需要到基层资源之间的距离；它们也不是一位改革者凭空造出的完整机器，而是在关中条件、既有行政实践和强邻压力中逐步展开。

这一重组提高了持续征发、补给与奖赏的能力，\eventref{WQ-0260-CHANGPING}{战国·长平之战}所需的长期竞争正可看见其重要性；但它也压缩了旧贵族的地方裁量，并让家庭更直接进入税、役、军功与相互监督的体系。这里的得失不是“集权较优”或“旧制较善”的裁断，而是国家能力的增长与社会承受的风险被同时放大。战国国家解决了春秋联盟难以常设化的问题，却把如何避免征发和惩罚失去限度的问题带入下一阶段。

\subsection{秦：统一的行政速度，与过度动员的回流}

\eventref{QIN-0230-UNIFICATION}{秦·统一六国}之后，郡县、文书、度量衡、道路与法律迅速进入新获土地。它们使不同地区能够以较一致的名目沟通、交换和征发，既结束了长期诸侯战争中的多重标准，也使远方开始接入中央的日常命令。统一并非因此成为没有摩擦的结果：地方官吏、仓储、道路和基层家庭必须不断承担行政运行所需的时间、物资与服役。

\eventref{QIN-0213-EMPIRE}{秦·高强度治理}所显示的风险，在于战国形成的战时动员没有随着统一自动降速。继承失序又使同一套命令网络暴露出脆弱处：\eventref{QIN-0210-SUCCESSION}{秦·沙丘继承}之后的中央不可信，与\eventref{QIN-0209-CHENSHENG}{秦·大泽起事}所引发的多点离散相互叠加，终于走向\eventref{QIN-0206-SURRENDER}{秦·子婴降汉}。这并不等于郡县、法律和统一尺度本身被历史否定；秦亡后它们仍被后继政权接续。真正留下的问题是，能够迅速集中资源的制度，如何在战争、工程与继承危机之外保有可承受的负担和可被信任的执行链。

\begin{turningpoint}[合卷结语：制度总在解决问题，也在分配代价]
西周以封国分担距离，春秋以霸主调度多国，战国以县邑和赏罚提高可计算性，秦再将这些能力推入统一帝国；四者面对的都不是抽象的“先进”或“落后”，而是不同规模下的治理、动员与整合。每次选择都把一部分成本转给地方精英、盟友、官吏或普通家庭，也把新的压力交给后来的政治行动者。读者可沿上述事件回到各章，观察这些压力怎样在具体的城邑、道路、会盟、战场与文书中发生作用。
\end{turningpoint}

\section{图表与地图索引}

\begin{description}
  \item[西周] \hyperref[fig:vol01:western-zhou:timeline]{“时间线：从牧野到东迁”}；\hyperref[fig:vol01:western-zhou:overview-map]{“空间总览：王畿、东方与边疆”}；\hyperref[fig:vol01:western-zhou:chengzhou]{“宗周、成周与东方网络示意”}；\hyperref[fig:vol01:western-zhou:zhao-southern-routes]{“昭王南向行动与汉水风险示意”}。
  \item[春秋] \hyperref[fig:vol01:spring-autumn:states-overview]{“空间总览：诸侯怎样围绕中原展开”}；\hyperref[fig:vol01:spring-autumn:xuge]{“繻葛之战空间示意”}；\hyperref[fig:vol01:spring-autumn:chengpu]{“城濮战役空间示意”}；\hyperref[fig:vol01:spring-autumn:boju]{“柏举战役空间示意”}；\hyperref[fig:vol01:spring-autumn:qin-jin-xiao-corridor]{“秦、晋、郑与崤函通道示意”}。
  \item[战国与秦] \hyperref[fig:vol01:warring-qin:seven-states]{“空间总览：七国竞争与秦的东进通道”}；\hyperref[fig:vol01:warring-qin:reform-network]{“战国变法关系示意”}；\hyperref[fig:vol01:warring-qin:changping]{“长平战场空间示意”}；\hyperref[fig:vol01:warring-qin:yan-qi]{“燕伐齐空间示意”}；\hyperref[fig:vol01:warring-qin:qin-unification]{“秦统一六国空间示意”}。
\end{description}
